\section{How to Move}

\begin{rulebox}{General Rule}
During his movement phase, the active player can move as many of his units and leaders as he desires in any direction up to the limit of each unit or leader's movement allowance, according to the procedures and within the restrictions detailed in this section. Artillery pieces cannot move. Except where specifically noted, the following movement restrictions apply to both units and leaders.
\end{rulebox}

\ruleslabel{Procedure}

Units and leaders move one at a time, tracing a path of connected hexes through the hex grid as they move. Each piece's movement allowance consists of a number of movement points (MPs). To enter a hex, the piece must spend one or more MPs depending on the terrain in that hex. In some cases, additional MPs must be spent in order to cross the terrain in the hexside through which the unit enters the hex. \textsc{Exception:} Disrupted units have special movement rules (11.8).

\caseheading{8.1} Only the active player can move his pieces during a movement phase.

He can move any of his pieces that he wants (other than artillery) in any order he wants. However, the active player cannot have combat during a movement phase. The inactive player cannot move or have combat during the active player's movement phase. Movement is always voluntary; a player is never forced to move.

\caseheading{8.2} Pieces must be moved individually.

Once a player moves a piece and withdraws his hand from that piece, the piece cannot be moved any further or have its move altered during that movement phase. In addition, once a player has begun to announce his attacks, he has irreversibly initiated his combat phase and cannot move any more pieces.

\caseheading{8.3} The presence of other pieces in a hex affects the ability of leaders and units to enter or to end their movement phase in that hex.

Units and leaders can always enter hexes occupied by other friendly units and leaders. There is no additional MP cost for entering or leaving a friendly-occupied hex. However, there can be only one unit (and any number of friendly leaders) in a hex at the end of a movement phase. Artillery pieces are not units and are unaffected by this rule. If there is more than one unit in a hex at the end of a movement phase, the active player must immediately eliminate (permanently remove from the game) all units in excess of one from that hex. Units removed in this fashion do count for army demoralization (13.0) and victory-point accumulation (15.0).

A piece can never enter a hex containing an enemy unit. If a unit enters a hex containing an enemy leader, the enemy leader is automatically placed in the nearest hex containing a friendly unit (of any force). There is no effect on the moving unit when it forces an enemy leader to displace to another hex. If there are two or more hexes into which a leader could be placed, the player controlling that leader decides to which hex the leader is displaced. Artillery pieces have no effect on movement.

\caseheading{8.4} Pieces can only enter connected hexes. Units and leaders cannot skip hexes during their movement.

\caseheading{8.5} A unit or leader cannot spend more MPs during a movement phase than it has in its movement allowance.

If a piece does not have sufficient MPs to pay all movement costs associated with entering a hex, it cannot enter that hex. A unit or leader does not have to spend all its available MPs; however, unused MPs cannot be saved for later use or loaned to another piece. Movement points that are not used during the movement phase in which they become available are lost. Unlike some games, \textsc{Cromwell's Victory} does not permit units or leaders to move a minimum of one hex regardless of their movement allowances.

\caseheading{8.6} The MP cost to enter each type of terrain on the map is listed in the Terrain Effects Chart.

It costs 1 MP to enter a clear hex. Other hexes can cost more to enter. To find the MP entry cost for a hex, find the type of terrain in that hex in the Type of Action column on the Terrain Effects Chart, then read across to the column representing the type of unit that is moving (Foot, Heavy Horse, or Light Horse). That number is the number of MPs it costs that unit type to enter that terrain type.

Leaders always pay 1 MP to enter all types of hex except town hexes (which no piece can enter, even via a lane), and they ignore crossing costs (see case 8.7). Their movement allowance is unaffected by visibility. Under certain circumstances, units and leaders can exit the map at a cost of 1 additional MP. See case 13.8 for details.

\caseheading{8.7} The MP cost to cross certain hexsides must be paid in addition to other MP costs to enter a hex.

Movement-point costs listed on the Terrain Effects Chart with a plus sign in front of them are hexside crossing costs that are added to the cost of entering a hex whenever that hex is entered across that type of hexside. If no extra cost is listed in a unit-type column, units of that type ignore that type of hexside for purposes of movement. Leaders always ignore all hexside terrain when moving.

\caseheading{8.8} Lanes negate the effects of other terrain for purposes of movement.

Whenever a unit enters a hex containing a lane via a hexside containing the lane symbol, the unit pays 1 MP to enter the hex, ignoring all hexside terrain that it can cross to enter the hex. Note, however, that a unit cannot enter a hex that is normally prohibited to it (contains a P on the Terrain Effects Chart) by moving along a lane. Lanes do not negate prohibitions against movement; they only allow units that could enter a hex anyway to do so at a lower cost in MPs. Lanes have no effect on leader movement or combat.

\caseheading{8.9} Enemy ZOCs (9.1) do not affect movement.
